\section{Prior art and claim boundary}\label{sec:prior-art}

\subsection{Homometry and reconstruction}

Patterson/autocorrelation data and homometric sets have a substantial
literature.  Rosenblatt and Seymour study structural aspects of homometric
sets \cite{RosenblattSeymour1982}; Mandereau et al. formulate left and right
interval contents in a group setting \cite{MandereauEtAl2011}; noncommutative
homometry in dihedral groups is treated by Genuys and Popoff
\cite{GenuysPopoff2018}.  Covariograms of nonconvex sets show related
reconstruction ambiguity in Euclidean geometry \cite{BenassiEtAl2011}.
Higher-order decks or bispectral data can carry information unavailable to a
second-order invariant \cite{RadcliffeScott2007,Kakarala2012}.

Our use of the right Patterson function is therefore classical at the level
of definition.  The finite result proved here is the exact collision-free
statement on the locked $82$ rows, together with the measured failure of
three label compressions.  We do not claim a general reconstruction theorem,
nor priority for group-ring homometry.  In particular, the 1989 paper of
Rosenblatt and Shapiro \cite{RosenblattShapiro1989} is cited as adjacent
semisimple-ring prior art; no theorem from it is imported.

\subsection{Coxeter complexes, coloured incidence, and tilings}

Coxeter complexes, weak order, and type-selected balanced complexes are
classical \cite{Bjorner1984a,Bjorner1984b}.  Coloured flag enumeration goes
back at least to Stanley's theory of balanced Cohen--Macaulay complexes and
its later refinements \cite{Stanley1979,Frohmader2012,White2022}.  Symbolic
chamber-system descriptions of tilings have likewise been developed in a
broader setting \cite{DressFriedrichsHuson1995,DressHuson1987}.

Exact factorisation of finite groups provides a natural algebraic language
for multiplier selectors \cite{BildanovEtAl2020}.  The typed carrier theorem
in \cref{thm:carrier} is not presented as a new general factorisation theory;
it is the complete finite structure of this particular $S_4$ atlas.

The word ``tile-transitive'' is also used narrowly.  Classical isohedral and
monotypic tiling theory concerns much broader ambient spaces and local-to-
global questions \cite{GrunbaumShephard1977,DolbilinSchulte2004}.  Our
partitions have one fixed tetrahedral container.  We neither produce a space
tiler nor classify tetrahedral space-fillers such as the families studied by
Sommerville and Goldberg \cite{Sommerville1924,Goldberg1974}.

\subsection{What is and is not classified}

The paper classifies connected chamber unions in one fixed $24$-chamber
$A_3$ subdivision, with the exact actions and cover notions of
\cref{sec:model}.  In a congruent Coxeter-pure cover, every member of the
partition remains a union of chambers of that fixed subdivision.  Its data
include every left-translate cover of each nonconvex tiling row and every
multiplier selector in those cover fibres.

It does not classify:
\begin{itemize}[leftmargin=2em]
\item disconnected chamber unions;
\item arbitrary polyhedral pieces or arbitrary tetrahedral dissections;
\item congruent copies whose boundaries cut interiors of the fixed chambers;
\item right-translate or two-sided group factorisations;
\item periodic or face-to-face tilings of $\mathbb R^3$;
\item Patterson reconstruction outside the displayed finite family;
\item contact up to unrestricted abstract complex isomorphism;
\item mechanical motion, hinging, or collision-free deformation;
\item a presentation-independent squared-edge invariant.
\end{itemize}

The count $15+82$ is already present in \cite{BabanskyyT2}.  The contribution
of the present work is the all-candidate congruence-rigidity theorem that
closes the earlier Coxeter-pure problem, together with the canonical
nonconvex classification and the exact structural layers that make the $82$
rows distinguishable and reviewable.
